% TOP_PROBLEM: {"id": 5500016, "title": "Simple Polygonalizations", "category_id": 6, "category": {"id": 6, "name": "geometry", "display_name": "Geometry", "description": "Euclidean and non-Euclidean geometry, geometric structures.", "slug": "geometry", "order_index": 6, "created_at": "2026-07-31T15:26:25.671Z"}, "status": "open", "problem_number": "AMR-054-0016", "set_id": 13}
% FIRST_POSTED: 2026-10-02
% ENTRY_KIND: statement_only
% UnsolvedMath ID 5500016; source catalogue code AMR-054-0016
% !TeX program = lualatex
% Definition and sources only; this file is not an LLM solution attempt.
\documentclass[11pt,a4paper]{article}
\usepackage[margin=25mm]{geometry}
\usepackage{fontspec}
\setmainfont{lmroman10-regular.otf}[BoldFont=lmroman10-bold.otf,
 ItalicFont=lmroman10-italic.otf,BoldItalicFont=lmroman10-bolditalic.otf]
\usepackage{amsmath,amssymb,mathtools}
\usepackage{unicode-math}
\setmathfont{latinmodern-math.otf}
\AtBeginDocument{\let\setminus\smallsetminus}
\usepackage{xurl,hyperref}
\hypersetup{colorlinks=true,urlcolor=blue}
\setcounter{secnumdepth}{0}
\setlength{\parindent}{0pt}
\setlength{\parskip}{5pt}
\setlength{\emergencystretch}{3em}
\begin{document}
\section*{Simple Polygonalizations}\label{5500016}
{\small UnsolvedMath ID 5500016 · AMR-054-0016 · Geometry · AMR Open Problem Lists}
\subsection{Definitions and mathematical statement}
Let $P\subset\mathbb R^2$ be a set of $n\ge3$ distinct points
in general position, meaning that no three are collinear.
A simple polygonalization is a cyclic ordering of all points
whose consecutive straight segments form a simple polygon:
nonadjacent edges are disjoint, and adjacent edges meet only
at their shared endpoint. Count two orderings as the same
polygonalization if one is a cyclic rotation or reversal of
the other. Write $\operatorname{poly}(P)$ for this count.

Can $\operatorname{poly}(P)$ be computed exactly in time polynomial
in $n$? In a bit model, take rational coordinates and require
running time polynomial in the total input bit length; in the
usual geometric model, elementary arithmetic and orientation
predicates are counted as constant-cost operations.

The output is the exact integer, not an estimate or a random
sample. Its binary length is polynomially bounded because there
are at most $(n-1)!/2$ candidate undirected cyclic orders.
Every point must be used exactly once as a polygon vertex;
additional Steiner vertices and polygonalizations of proper
subsets are excluded. Counting oriented or rooted cyclic orders
changes the answer by a known factor and does not alter the
complexity question. Efficient algorithms for restricted point
sets do not settle the question for arbitrary planar inputs.
\subsection{Short English statement}
Can the exact number of simple polygonalizations of a planar point set be computed in polynomial time?
\subsection{Sources}
\begin{itemize}
\item[S1] ULAM.AI and the UnsolvedMath Contributors, supplied problems.json, record 5500016 (AMR-054-0016), AMR Open Problem Lists. Catalogue identity and source transcription; CC BY 4.0.
\url{https://huggingface.co/datasets/ulamai/UnsolvedMath}.
\item[S2] Open Problems Project - Computational Geometry, Problem 16. Original source indexed by AMR-054-0016.
\url{https://topp.openproblem.net/p16}.
\item[S3] The Open Problems Project, original numbered problem and status notes.
\url{https://topp.openproblem.net/p16}.
\end{itemize}
\end{document}
